\clearpage\thispagestyle{empty}
\pdfbookmark[0]{Quellen zum Bericht über den Merkurdurchgang von 1802}{mercur18021103-sources}
\begin{center}\textbf{Quellen zum Bericht über den Merkurdurchgang von 1802}\end{center}
\noindent Diese KI-Übertragung in heutiges Deutsch folgt dem Auszug aus Franz Xaver von Zachs Bericht über Gauß’ Beobachtungen in Braunschweig in \href{https://gdz.sub.uni-goettingen.de/id/PPN236007467}{Gauß’ \textit{Werken}, Band~VI} (Göttingen, 1874), S.~231--232. Der Auszug enthält die Zeiten beider Berührungen beim Austritt, die mittleren Zeiten in einem Handschriftenauszug in eckigen Klammern und fünf Beobachtungen der Rektaszensionsdifferenz zwischen Sonne und Merkur.
\par\smallskip\noindent Der \href{https://api.digitale-sammlungen.de/iiif/image/v2/bsb10538598_00581/full/1450,/0/default.jpg}{Erstdruck} steht als Beitrag LXIII in der \textit{Monatlichen Correspondenz}, Band~VI, S.~567--575. Der Auszug reicht von \href{https://api.digitale-sammlungen.de/iiif/image/v2/bsb10538598_00584/full/1450,/0/default.jpg}{Gauß’ Beobachtungen auf S.~570} bis zu \href{https://api.digitale-sammlungen.de/iiif/image/v2/bsb10538598_00585/full/1450,/0/default.jpg}{seiner Tabelle auf S.~571}. Die Überschrift in den \textit{Werken} nennt den 3. November 1802. Die Zeitschrift nennt in der Überschrift den 8. November, der Handschriftenauszug in den \textit{Werken} den 9. November. Die Überschrift und das Datum der Handschrift folgen den \textit{Werken}.
\par\smallskip\noindent Übertragung und Satz wurden mit KI erstellt. Die Digitalisate der Zeitschrift stammen von der Bayerischen Staatsbibliothek.
\clearpage
